% =============================================================================
% DOSSIÊ DE ESTUDO — MOLAMBUDOS: O DIÁRIO DO PACIENTE 1.260
% Aparato crítico. NÃO é parte da ficção e não deve ser encadernado com o miolo.
% =============================================================================
\documentclass[11pt,a4paper]{article}

\usepackage[T1]{fontenc}
\usepackage[utf8]{inputenc}
\usepackage[brazilian]{babel}
\usepackage{ebgaramond}
\usepackage[margin=2.6cm]{geometry}
\usepackage{xcolor}
\usepackage{booktabs}
\usepackage{longtable}
\usepackage{enumitem}
\usepackage{titlesec}
\usepackage[hidelinks]{hyperref}

\definecolor{sangue}{HTML}{8B0000}
\definecolor{tinta}{HTML}{2E241A}
\titleformat{\section}{\Large\bfseries\color{sangue}}{\thesection}{0.7em}{}
\titleformat{\subsection}{\large\bfseries\color{tinta}}{\thesubsection}{0.7em}{}
\setlist{itemsep=0.25em, topsep=0.35em}

\newcommand{\ressalva}[1]{%
  \par\vspace{0.5em}\noindent\fcolorbox{sangue!40}{sangue!4}{%
  \parbox{0.965\linewidth}{\small\textbf{Ressalva metodológica.} #1}}\par\vspace{0.5em}}

\title{\vspace{-1.5cm}\textbf{Molambudos --- O Diário do Paciente 1.260}\\[0.3em]
\large Dossiê de estudo: projeto literário, técnicas, teorias, limites}
\author{\small Aparato crítico da edição trilíngue}
\date{\small 5 de agosto de 2026}

\begin{document}
\maketitle
\thispagestyle{empty}

\begin{abstract}
\noindent Este dossiê descreve e analisa \textit{Molambudos --- O Diário do
Paciente 1.260}, romance-arquivo em hipertexto impresso, publicado em três
edições de idioma e uma trilíngue. Reúne a ficha técnica medida da obra, sua
arquitetura formal, a linhagem literária em que se inscreve, as técnicas
narrativas que emprega, os fundamentos teóricos que mobiliza, sua base
histórica, e uma análise crítica que inclui limitações verificadas. É um
aparato de estudo, não uma peça promocional: as seções \ref{sec:limites} e
\ref{sec:validacao} existem para delimitar o que a obra \textit{não} alcançou
e o que ainda \textit{não} foi validado externamente.
\end{abstract}

\ressalva{Todos os números deste dossiê foram medidos sobre os arquivos-fonte
da obra em 2026-08-05, não estimados. Onde há inferência interpretativa, o
texto diz explicitamente que infere. Onde a afirmação exigiria validação
externa inexistente --- recepção, impacto, comparação de mérito ---, o dossiê
se abstém em vez de estimar. Ver \S\ref{sec:validacao}.}

% =============================================================================
\section{Ficha técnica}

\begin{center}\small
\begin{tabular}{@{}ll@{}}
\toprule
\textbf{Título} & Molambudos --- O Diário do Paciente 1.260 \\
\textbf{Autoria} & Marcelo Dias de Carvalho Filho \\
\textbf{Coautoria e ilustração} & Marcelo Claro \\
\textbf{ISBN} & 9798189170492 \\
\textbf{Gênero} & Romance-arquivo; horror psicológico documental \\
\textbf{Forma} & Hipertexto impresso, 84 fragmentos, 216 rotas de leitura \\
\textbf{Idiomas} & Português (PT-BR), inglês, chinês simplificado \\
\textbf{Formato de impressão} & 160 $\times$ 230 mm, brochura \\
\bottomrule
\end{tabular}
\end{center}

\vspace{0.6em}
\noindent\textbf{Extensão medida por edição}

\begin{center}\small
\begin{tabular}{@{}lrrr@{}}
\toprule
\textbf{Edição} & \textbf{Miolo digital} & \textbf{Impressão 160$\times$230} & \textbf{Palavras} \\
\midrule
Português   & 415 p. & 435 p. & 43.730 \\
Inglês      & 411 p. & 427 p. & 44.978 \\
Chinês      & 387 p. & 397 p. & 70.750 caracteres \\
Trilíngue   & 1.061 p. & 1.115 p. & --- \\
\bottomrule
\end{tabular}
\end{center}

\noindent\small A contagem em palavras não é comparável entre o chinês e as
línguas ocidentais; para o chinês registra-se a contagem de caracteres, que é
a métrica usual.

% =============================================================================
\section{O projeto literário}

\textit{Molambudos} apresenta-se como um \textbf{arquivo pericial}, não como
um romance. O leitor recebe 84 fragmentos --- memórias, laudos, cartas,
transcrições de áudio, formulários clínicos --- reunidos por um arquivista
anônimo em torno de um objeto: o diário de Joaquim Antônio Correia,
registrado no Hospital Colônia de Barbacena sob o número 1.260.

A proposta declarada, legível na estrutura e no paratexto, articula três
movimentos:

\begin{enumerate}
\item \textbf{Restituir uma história soterrada.} O Colônia de Barbacena
existiu; nele morreram, ao longo do século XX, um número de pessoas que a
historiografia estima em dezenas de milhares, em boa parte sem qualquer
diagnóstico psiquiátrico. A obra constrói sobre esse solo.

\item \textbf{Fazer da forma um argumento.} A fragmentação não é ornamento:
ela encena a impossibilidade de reconstituir uma vida a partir de restos
documentais. O leitor tem de montar o caso, e a montagem é sempre parcial.

\item \textbf{Implicar o leitor.} A obra recusa a posição de espectador. Uma
família inteira de fragmentos (CONT) endereça o leitor em segunda pessoa e o
inscreve na cadeia de transmissão como \textbf{Paciente 1.263}.
\end{enumerate}

\subsection{A cadeia de transmissão}

O dispositivo central é uma fila numerada de corpos ocupados:

\begin{center}\small
\begin{tabular}{@{}cll@{}}
\toprule
\textbf{Nº} & \textbf{Quem} & \textbf{Onde se fixa} \\
\midrule
1.259 & o anterior a Joaquim & CONT-07, CONT-08 \\
1.260 & Joaquim Antônio Correia & título da obra \\
1.261 & Dr.\ Heitor Oliveira & DOC-16 (assina) \\
1.262 & Dra.\ Lúcia Mendes & LUC-10, LUC-11, Epílogo \\
\textbf{1.263} & \textbf{o leitor} & LUC-12, DOC-24, CONT-09, protocolo \\
\bottomrule
\end{tabular}
\end{center}

A precisão dessa contagem é condição de funcionamento: é ela que transforma
a última posição em ameaça pessoal em vez de metáfora vaga.

% =============================================================================
\section{Arquitetura formal}

\subsection{As quatro famílias de fragmento}

\begin{center}\small
\begin{tabular}{@{}llrrr@{}}
\toprule
\textbf{Sigla} & \textbf{Natureza} & \textbf{N} & \textbf{Palavras} & \textbf{Média} \\
\midrule
MEM  & Memórias do Paciente 1.260 & 27 & 13.310 & 492 \\
DOC  & Documentos, laudos, cartas & 27 & 15.247 & 564 \\
LUC  & Investigação de Lúcia Mendes (2026) & 17 & 8.934 & 525 \\
CONT & Protocolos de contaminação do leitor & 13 & 6.239 & 479 \\
\midrule
     & \textbf{Total} & \textbf{84} & \textbf{43.730} & 521 \\
\bottomrule
\end{tabular}
\end{center}

A variação de extensão é acentuada e funcional: de 130 palavras (MEM-26) a
2.310 (DOC-03, o laudo de necropsia). Os fragmentos curtos operam como golpes;
os longos, como imersão documental.

\subsection{Cinco partes, três atos}

A leitura linear percorre cinco partes --- Sertão (9), Colônia (16), Diário de
Oliveira e Laudos (25), Investigação Lúcia (20), Contaminação (14 + Epílogo)
--- que se agrupam em três atos: o Trauma (9 fragmentos), a
Institucionalização (41) e o Ciclo (35).

A assimetria é significativa: o ato central, que corresponde aos 62 anos de
internação, ocupa quase metade da obra. A duração do confinamento é
reproduzida na duração da leitura.

\subsection{A rede de rotas}

\begin{center}\small
\begin{tabular}{@{}lr@{}}
\toprule
Fragmentos & 84 \\
Rotas de leitura (\texttt{\textbackslash rota}) & 216 \\
Destinos distintos & 85 (84 fragmentos + Epílogo) \\
Fragmentos sem rota de entrada & 0 \\
Fragmentos de que se alcança o Epílogo & 84 / 84 \\
Ciclos (componentes fortemente conexos $>$ 1) & 6 \\
Ciclos sem saída & 0 \\
Nó mais referenciado & DOC-01, com 20 rotas de entrada \\
\bottomrule
\end{tabular}
\end{center}

Duas propriedades são deliberadas e foram verificadas: \textbf{toda rota leva
ao Epílogo} e \textbf{nenhum ciclo aprisiona o leitor}. Os ciclos são
preservados --- são a tese da obra, que afirma ``o ciclo recomeça'' --- mas
nenhum é uma armadilha: de qualquer fragmento existe percurso finito até o
fim.

\ressalva{A alternativa seria um grafo acíclico estrito. Ela foi medida e
descartada: custaria 85 das 216 rotas, precisamente as remissões
documento$\to$memória que dão à obra o caráter de arquivo, e contradiria a
tese do texto. A escolha é do autor e está registrada.}

% =============================================================================
\section{Linhagem literária e narrativa}

\ressalva{Esta seção identifica \textbf{afinidades formais verificáveis no
texto}. Não afirma que o autor leu as obras citadas nem que delas derivou
conscientemente. Linhagem, aqui, é parentesco de procedimento --- o que a
crítica pode comparar ---, não genealogia de leitura.}

\subsection{O romance-arquivo e o documento forjado}

A obra inscreve-se na tradição do \textit{found footage} literário, cuja
matriz moderna é epistolar: \textit{Drácula} (1897) monta o horror a partir de
diários, cartas e recortes, e delega ao leitor a síntese que nenhum
personagem alcança. \textit{Molambudos} adota o mesmo princípio e o radicaliza
ao incluir laudos periciais, cadeias de custódia e pareceres de conservação
--- gêneros burocráticos, não íntimos.

Aproxima-se também de \textit{Casa de Folhas} (Danielewski, 2000) no uso da
página como espaço, dos aparatos de navegação e da vertigem bibliográfica.
Difere num ponto decisivo: onde Danielewski multiplica níveis de comentário
até o labirinto, \textit{Molambudos} mantém a cadeia curta e nominal, e a usa
para chegar ao leitor.

\subsection{A ficção hipertextual em suporte impresso}

\textit{O Jogo da Amarelinha} (Cortázar, 1963) estabeleceu o precedente do
livro com mais de uma ordem de leitura. \textit{Molambudos} vai além do
``tabuleiro'' cortazariano ao oferecer uma rede densa (216 arestas entre 85
nós) com propriedades garantidas por verificação automatizada --- convergência
e ausência de aprisionamento.

\ressalva{Essa verificação é uma característica de \textit{produção}, não de
mérito literário. Ela garante que a promessa de navegação seja cumprida; não
diz nada sobre a qualidade da prosa.}

\subsection{O testemunho e a literatura de atrocidade}

A obra dialoga com a tradição brasileira de denúncia do sistema manicomial,
cujo marco jornalístico contemporâneo é \textit{Holocausto Brasileiro}
(Daniela Arbex, 2013), citada nominalmente em DOC-19. Em registro literário,
ecoa \textit{Vidas Secas} (Graciliano Ramos, 1938) na primeira parte: a seca,
a marcha da família, a linguagem despojada.

A questão teórica de fundo é a da \textit{representabilidade}: como narrar uma
atrocidade sem espetacularizá-la. A escolha da obra --- horror sobrenatural
como metáfora explícita, declarada como tal na \textit{Nota Histórica} --- é
uma posição, e uma posição discutível, que a \S\ref{sec:critica} examina.

\subsection{A segunda pessoa e o leitor implicado}

O endereçamento em segunda pessoa tem precedente canônico em \textit{Se um
viajante numa noite de inverno} (Calvino, 1979), onde é lúdico e
metaliterário. \textit{Molambudos} o emprega em chave inversa: não para
brincar com a leitura, mas para acusá-la. Vinte e sete dos 84 fragmentos
mobilizam o ``você'' de forma sustentada.

% =============================================================================
\section{Técnicas narrativas}

\subsection{Indução e sugestão}

Os fragmentos CONT empregam procedimentos identificáveis da literatura de
sugestão clínica:

\begin{itemize}
\item \textbf{Duplo vínculo.} CONT-03 constrói uma alternativa sem saída:
``\textit{Se sim: o cheiro está lá. [\ldots] Se não: mentira. Você levantou.
Todos levantam.}'' Ambas as respostas confirmam a premissa.

\item \textbf{Sugestão ideomotora.} ``\textit{Mais tarde, quando você estiver
distraído, seu dedo vai voltar à página sozinho.}''

\item \textbf{Sugestão pós-hipnótica.} ``\textit{Ele vai estar no seu
travesseiro esta noite.}''

\item \textbf{Ancoragem fisiológica.} A descrição de sintomas (respiração,
pupilas, temperatura das extremidades) tende a induzir a atenção que os
constata.

\item \textbf{Pareidolia dirigida.} ``\textit{Olhe para o canto superior
direito do quarto [\ldots] Fixe o olho no canto escuro por 10 segundos.}''
\end{itemize}

\subsection{A inversão da ausência}

A técnica mais sofisticada da obra está em CONT-04, ``Noite 4'': após três
noites de pesadelo, nada acontece. E o texto converte o alívio em ameaça:
``\textit{a ausência do medo era pior que o medo [\ldots] A espera é o
verdadeiro contágio.}''

O procedimento é notável porque remove o estímulo e mantém o efeito ---
transfere a produção do horror do texto para o leitor.

\subsection{O instrumento clínico como forma}

DOC-09 e CONT-05 adotam a forma de escalas de autoavaliação psiquiátrica,
com itens, alternativas e escore final. O item 15 de DOC-09 --- ``\textit{Se
você pudesse parar de ler agora e nunca mais abrir este livro, você pararia?
( ) Sim ( ) Não ( ) Já tentei}'' --- opera por duplo sentido: dentro da ficção
refere-se à leitura; no registro do instrumento clínico que imita, lê-se como
item de suicidalidade.

\ressalva{Essa ambiguidade é eficaz e é também a razão pela qual a página de
aviso inicial, com a linha do CVV, não é ornamento. Ver \S\ref{sec:etica}.}

\subsection{Marginália e paratexto diegético}

A obra emprega notas marginais (\texttt{margnota}) que comentam o texto de
fora, fichas de paciente, selos de arquivo e uma página de ``Cuidado''
integralmente diegética. O paratexto participa da ficção em vez de emoldurá-la.

% =============================================================================
\section{Fundamentos teóricos mobilizados}

\ressalva{A obra não declara filiação teórica. O que segue são os quadros que
um estudo crítico pode aplicar produtivamente, com indicação do que na obra
os convoca --- não atribuição de intenção ao autor.}

\begin{description}[style=unboxed,leftmargin=0pt]
\item[Trauma e transmissão intergeracional.] A criatura funciona como
figuração da dor que atravessa gerações sem ser elaborada. A própria
\textit{Nota Histórica} explicita: ``\textit{a dor que atravessa gerações, que
ninguém processa e que passa adiante como um contágio, é a única criatura
verdadeira}''. É a chave de leitura que a obra oferece de si.

\item[Instituição total.] O Colônia é descrito com os traços que Goffman
identificou nas instituições totais: ruptura das barreiras entre esferas da
vida, mortificação do eu, economia de privilégios. A ala dos silenciados
opera como intensificação desse regime.

\item[Necropolítica e vidas descartáveis.] A obra insiste em que os mortos do
Colônia eram ``pobres, negras, epilépticas, mães solteiras, homossexuais''. O
critério de internação é social, não clínico --- o que desloca a leitura do
horror psiquiátrico para o horror político.

\item[Arquivo, poder e lacuna.] A cadeia de custódia com lacunas documentais
(DOC-20) e as ocorrências de arquivo (DOC-23) tematizam o arquivo como
instância de poder e o silêncio documental como violência adicional.

\item[Recepção e leitor implicado.] A obra constrói um leitor-modelo e o
nomeia. Toda a família CONT é um dispositivo de teoria da recepção posto a
funcionar em vez de discutido.
\end{description}

% =============================================================================
\section{Base histórica e postura ética}
\label{sec:etica}

A \textit{Nota Histórica} da obra é metodologicamente cuidadosa, e convém
registrá-lo porque é raro no gênero:

\begin{itemize}
\item Declara que Joaquim, a criatura e a instituição narrada são ficcionais.
\item Distingue dois referentes históricos que a ficção poderia confundir: o
curral de retirantes de 1915 (Alagadiço) e o de Senador Pompeu (1932--1933).
\item Afirma explicitamente que ``\textit{os totais de internados e mortos não
são apresentados como precisos}''.
\item Atribui a base documental do Colônia de Barbacena a fonte nomeada
(Arbex).
\item Declara o horror sobrenatural como metáfora, não como tese.
\end{itemize}

O volume traz ainda página de aviso ao leitor, anterior à ficção, com
enumeração de conteúdos sensíveis, distinção ficção/história e telefone do
Centro de Valorização da Vida. Trata-se de paratexto: precede o quadro
ficcional e não interrompe a imersão.

% =============================================================================
\section{Análise crítica}
\label{sec:critica}

\subsection{Onde a obra é forte}

\begin{enumerate}
\item \textbf{A forma sustenta o argumento.} A fragmentação, a assimetria
entre as partes e a rede de rotas não são decorativas: reproduzem
formalmente a impossibilidade de reconstituição e a duração do confinamento.

\item \textbf{A implicação do leitor é executada, não anunciada.} A cadeia
numerada termina numa posição vaga que o leitor ocupa. É um dispositivo
simples e rigorosamente construído.

\item \textbf{A honestidade documental é explícita.} A obra separa o que
inventa do que documenta, e recusa precisão que não tem.

\item \textbf{A prosa de indução é tecnicamente competente.} A inversão de
CONT-04 é o melhor exemplo: retira o estímulo e conserva o efeito.
\end{enumerate}

\subsection{Onde a obra é discutível}

\begin{enumerate}
\item \textbf{A metáfora sobrenatural pode aliviar o horror histórico.} Ao
dar ao sofrimento uma causa fantástica, a obra corre o risco de deslocar a
responsabilidade política para uma entidade. A \textit{Nota Histórica}
antecipa a objeção e a responde; se a resposta basta é questão crítica aberta,
e é a mais produtiva que a obra oferece ao debate.

\item \textbf{As notas do editor interrompem a indução.} Treze fragmentos
trazem notas de rodapé que declaram o dispositivo como formal. Duas caem no
pico do efeito --- CONT-03 encerra em ``\textit{Para seu corpo, ela é real}'' e
a nota emenda informando que a contaminação é textual. A informação é
necessária; a posição é discutível. Consolidá-las em nota única ao final
preservaria ambas as funções.

\item \textbf{A densidade de segunda pessoa pode saturar.} Vinte e sete
fragmentos mobilizam o ``você''. Em leitura linear contínua, o procedimento
tende ao previsível --- risco que a leitura por rotas mitiga, mas que a
edição impressa não elimina.
\end{enumerate}

% =============================================================================
\section{Limitações verificadas}
\label{sec:limites}

\begin{enumerate}
\item \textbf{As três edições não são o mesmo livro.} Dez fragmentos divergem
estruturalmente além do que a diferença de idioma explica. Os casos extremos:
CONT-05 tem 163 linhas em PT e EN e 63 em ZH; CONT-03 tem 174 em PT contra 102
em EN e ZH; DOC-26 tem 110 em PT e EN contra 56 em ZH. Na edição trilíngue as
três versões aparecem lado a lado, o que torna a divergência visível ao
leitor. \textbf{É a limitação mais séria para a pretensão de obra de
referência internacional.}

\item \textbf{A arte de capa está abaixo do padrão de impressão.} Os arquivos
de capa e contracapa têm 1.046 $\times$ 1.504 px, o que resulta em 163 DPI no
painel impresso; o padrão exigido é 300 DPI.

\item \textbf{A contracapa traz um código de barras de teste.} O elemento
embutido na arte declara ISBN 978-65-01-23456-7, que não é o ISBN da obra.

\item \textbf{O corpus não está sob controle de versão.} Os arquivos-fonte da
obra não têm histórico de commits: não há diff nem possibilidade de reverter
uma edição malfeita.
\end{enumerate}

% =============================================================================
\section{Estado de validação}
\label{sec:validacao}

Esta seção existe para impedir que o dossiê seja lido como aval.

\begin{center}\small
\begin{tabular}{@{}p{6.5cm}l@{}}
\toprule
\textbf{Dimensão} & \textbf{Estado} \\
\midrule
Integridade de compilação das 5 edições & verificada \\
Paridade de rotas fonte $\times$ impresso & verificada (648/648) \\
Convergência e ausência de aprisionamento & verificada \\
Integridade dos fragmentos (SHA-256) & verificada \\
\midrule
Revisão por pares acadêmica & \textbf{não realizada} \\
Revisão editorial profissional externa & \textbf{não realizada} \\
Verificação historiográfica por especialista & \textbf{não realizada} \\
Revisão nativa das traduções EN e ZH & \textbf{não realizada} \\
Recepção crítica & \textbf{inexistente} (obra não publicada) \\
Impacto & \textbf{não mensurável} \\
\bottomrule
\end{tabular}
\end{center}

\ressalva{As quatro primeiras linhas são verificações \textbf{técnicas
internas}: dizem que o objeto está bem construído, não que é bom. Nenhuma
delas substitui julgamento crítico. As seis últimas não foram realizadas, e
este dossiê não as antecipa. Qualquer afirmação de mérito, originalidade ou
impacto que apareça em material derivado deste texto deve ser tratada como
não fundamentada até que exista validação externa.}

% =============================================================================
\section{Agenda de pesquisa}

Questões que a obra abre e que um estudo poderia enfrentar:

\begin{enumerate}
\item A metáfora sobrenatural amplia ou atenua a força testemunhal de uma
narrativa sobre atrocidade documentada?
\item Que efeitos a rede de rotas produz na experiência de leitura, medidos
empiricamente, e não postulados?
\item O endereçamento em segunda pessoa sustentado por 27 fragmentos produz
implicação ou reatância no leitor?
\item Como a divergência entre as três edições altera a obra recebida em cada
língua --- e o que isso diz sobre o próprio conceito de ``edição trilíngue''?
\item De que modo os paratextos de proteção (aviso, notas do editor)
interagem com o dispositivo de imersão que a obra constrói?
\end{enumerate}

% =============================================================================
\section*{Referências mencionadas}
\addcontentsline{toc}{section}{Referências mencionadas}

\begin{itemize}\small
\item ARBEX, Daniela. \textit{Holocausto Brasileiro}. Geração Editorial, 2013.
\item CALVINO, Italo. \textit{Se um viajante numa noite de inverno}. 1979.
\item CORTÁZAR, Julio. \textit{O Jogo da Amarelinha}. 1963.
\item DANIELEWSKI, Mark Z. \textit{House of Leaves}. Pantheon, 2000.
\item GOFFMAN, Erving. \textit{Manicômios, prisões e conventos}. 1961.
\item RAMOS, Graciliano. \textit{Vidas Secas}. 1938.
\item STOKER, Bram. \textit{Drácula}. 1897.
\end{itemize}

\vfill
\noindent\rule{\linewidth}{0.4pt}\\
{\footnotesize Dossiê gerado como aparato crítico. Sobre ressalvas de métricas
e alegações do projeto, ver \texttt{CORRIGENDUM.md}.}

\end{document}
